\documentclass[../main.tex]{subfiles}
\begin{document}

\chapter{Inter-Process Communication}
\lessoninfo{
  video={P3L3},
  textbook={Silberschatz et al.\ \cite{silberschatz2018} ch.~3; OSTEP \cite{ostep} ch.~5; Kerrisk \cite{kerrisk2010} ch.~43--54, 56},
  keywords={inter-process communication, message passing, port, pipe, message queue, socket, shared memory, copy versus map, local procedure call, shared memory segment, System V IPC, POSIX shared memory, memory-mapped file, process-shared mutex},
}

Lesson~2 introduced inter-process communication (IPC) as the way in which
processes, each isolated in its own address space, exchange data, and it
distinguished message passing from shared memory.  This lesson examines the
mechanisms that operating systems provide for both.  Shared memory builds on
the address translation of Lesson~8, which lets the operating system map the
same physical memory into several address spaces; it keeps the kernel out of
every exchange, but it leaves the synchronization and the communication
protocol, which the kernel handles for message passing, to the processes.

\section{Inter-process communication}
\label{sec:l09-ipc}

IPC comprises the mechanisms that the operating system supports for the
interaction of processes: \emph{communication}, the transfer of data from
one address space to another, and \emph{coordination} of the processes'
actions.\source{P3L3, inter-process communication; Silberschatz et al.\
ch.~3.}\index{inter-process communication}  The mechanisms fall into four
groups.
\begin{description}
  \item[Message passing] Processes exchange messages through a channel that
    the kernel maintains, e.g.\ sockets, pipes and message queues.
  \item[Memory-based IPC] The operating system makes memory accessible to
    several processes, which communicate by reading and writing it, e.g.\
    shared memory and memory-mapped files.
  \item[Higher-level semantics] Mechanisms that give the data transfer more
    structure: files that several processes read and write, and remote
    procedure calls (Lesson~13).
  \item[Synchronization primitives] Constructs such as mutexes and
    semaphores with which processes coordinate (Lesson~10).
\end{description}
This lesson concentrates on message passing, on shared memory, and on the
synchronization that shared memory requires.

\section{Message-based IPC}
\label{sec:l09-message}

In message-based IPC\index{message passing} the processes create messages
and send and receive them through a channel that the operating system
creates and maintains, typically a buffer in kernel memory organized as a
FIFO queue.\source{P3L3, message-based IPC; Silberschatz et al.\ ch.~3.}
The processes never access this buffer directly.

\begin{definition}[Port]
  A \term{port}[ポート] is the interface through which a process uses a
  message-passing channel: the process sends (writes) messages to a port and
  receives (reads) messages from a port, and the kernel moves the messages
  between the ports.
\end{definition}

The kernel is involved both in establishing the communication and in every
single IPC operation.  A send is a system call that copies the message from
the address space of the sender into the channel; a receive is a system call
that copies it from the channel into the address space of the receiver.  A
request-response exchange between a client and a server therefore costs
four user/kernel crossings and four data copies (\cref{fig:l09-message-ipc}).

\begin{figure}[htbp]
  \centering
  % Message-based IPC: a request-response exchange through a kernel channel.
% Every numbered crossing is one system call and one data copy.
% Usage: % Message-based IPC: a request-response exchange through a kernel channel.
% Every numbered crossing is one system call and one data copy.
% Usage: % Message-based IPC: a request-response exchange through a kernel channel.
% Every numbered crossing is one system call and one data copy.
% Usage: \input{figures/l09-message-ipc}   (width ~116 mm)
\begin{tikzpicture}[x=1mm, y=1mm, font=\footnotesize\sffamily,
  proc/.style={draw, rounded corners=2pt, fill=ln-box, align=center,
               minimum width=26mm, minimum height=9mm, inner sep=1pt},
  port/.style={draw, fill=ln-accent!15, minimum width=4mm,
               minimum height=21mm, inner sep=0pt},
  cell/.style={draw, fill=white, minimum width=4mm, minimum height=4.5mm,
               inner sep=0pt},
  msgc/.style={cell, fill=ln-accent!35},
  stepn/.style={circle, draw, fill=white, inner sep=0.5pt,
               minimum size=3.6mm, font=\scriptsize\sffamily\bfseries},
  lab/.style={font=\scriptsize\sffamily, text=ln-muted},
  >={Stealth[length=1.6mm]}]
  % processes
  \node[proc] (p1) at (14,28) {P1 (\iflangja{クライアント}{client})};
  \node[proc] (p2) at (98,28) {P2 (\iflangja{サーバ}{server})};
  % user/kernel boundary
  \draw[dashed, ln-rule, thick] (0,19) -- (116,19);
  \node[lab, anchor=south] at (56,19.5) {\iflangja{ユーザモード}{user mode}};
  \node[lab, anchor=north] at (56,18.5) {\iflangja{カーネルモード}{kernel mode}};
  % channel
  \draw[rounded corners=2pt, fill=ln-box] (26,-20) rectangle (86,10);
  \node[port] (pt1) at (32,-3.5) {};
  \node[port] (pt2) at (80,-3.5) {};
  \node[lab, anchor=north] at (32,-14.6) {\iflangja{ポート}{port}};
  \node[lab, anchor=north] at (80,-14.6) {\iflangja{ポート}{port}};
  \foreach \i in {0,...,7} {
    \pgfmathsetmacro{\x}{42+4*\i}
    \ifnum\i>5 \node[msgc] at (\x,1.5) {}; \else \node[cell] at (\x,1.5) {}; \fi
    \ifnum\i<2 \node[msgc] at (\x,-8.5) {}; \else \node[cell] at (\x,-8.5) {}; \fi
  }
  \draw[->] (34,1.5) -- (40,1.5);
  \draw[->] (72,1.5) -- (78,1.5);
  \draw[->] (78,-8.5) -- (72,-8.5);
  \draw[->] (40,-8.5) -- (34,-8.5);
  \node[lab, anchor=south] at (56,4.2) {\iflangja{要求}{requests}};
  \node[lab, anchor=north] at (56,-11.2) {\iflangja{応答}{responses}};
  \node[lab, anchor=north] at (56,-20.5)
    {\iflangja{チャネル（カーネル内のバッファ）}{channel (kernel buffers)}};
  % the four system calls
  \draw[->, thick] (22,23.5) -- (22,1.5) -- (30,1.5);
  \draw[->, thick] (82,1.5) -- (90,1.5) -- (90,23.5);
  \draw[->, thick] (106,23.5) -- (106,-8.5) -- (82,-8.5);
  \draw[->, thick] (30,-8.5) -- (6,-8.5) -- (6,23.5);
  \node[stepn] at (22,19) {1};
  \node[stepn] at (90,19) {2};
  \node[stepn] at (106,19) {3};
  \node[stepn] at (6,19) {4};
  \node[lab, anchor=west, text=black] at (24.2,13) {\texttt{send}};
  \node[lab, anchor=east, text=black] at (88,13)   {\texttt{recv}};
  \node[lab, anchor=east, text=black] at (104,13)  {\texttt{send}};
  \node[lab, anchor=west, text=black] at (8.2,13)  {\texttt{recv}};
\end{tikzpicture}
   (width ~116 mm)
\begin{tikzpicture}[x=1mm, y=1mm, font=\footnotesize\sffamily,
  proc/.style={draw, rounded corners=2pt, fill=ln-box, align=center,
               minimum width=26mm, minimum height=9mm, inner sep=1pt},
  port/.style={draw, fill=ln-accent!15, minimum width=4mm,
               minimum height=21mm, inner sep=0pt},
  cell/.style={draw, fill=white, minimum width=4mm, minimum height=4.5mm,
               inner sep=0pt},
  msgc/.style={cell, fill=ln-accent!35},
  stepn/.style={circle, draw, fill=white, inner sep=0.5pt,
               minimum size=3.6mm, font=\scriptsize\sffamily\bfseries},
  lab/.style={font=\scriptsize\sffamily, text=ln-muted},
  >={Stealth[length=1.6mm]}]
  % processes
  \node[proc] (p1) at (14,28) {P1 (\iflangja{クライアント}{client})};
  \node[proc] (p2) at (98,28) {P2 (\iflangja{サーバ}{server})};
  % user/kernel boundary
  \draw[dashed, ln-rule, thick] (0,19) -- (116,19);
  \node[lab, anchor=south] at (56,19.5) {\iflangja{ユーザモード}{user mode}};
  \node[lab, anchor=north] at (56,18.5) {\iflangja{カーネルモード}{kernel mode}};
  % channel
  \draw[rounded corners=2pt, fill=ln-box] (26,-20) rectangle (86,10);
  \node[port] (pt1) at (32,-3.5) {};
  \node[port] (pt2) at (80,-3.5) {};
  \node[lab, anchor=north] at (32,-14.6) {\iflangja{ポート}{port}};
  \node[lab, anchor=north] at (80,-14.6) {\iflangja{ポート}{port}};
  \foreach \i in {0,...,7} {
    \pgfmathsetmacro{\x}{42+4*\i}
    \ifnum\i>5 \node[msgc] at (\x,1.5) {}; \else \node[cell] at (\x,1.5) {}; \fi
    \ifnum\i<2 \node[msgc] at (\x,-8.5) {}; \else \node[cell] at (\x,-8.5) {}; \fi
  }
  \draw[->] (34,1.5) -- (40,1.5);
  \draw[->] (72,1.5) -- (78,1.5);
  \draw[->] (78,-8.5) -- (72,-8.5);
  \draw[->] (40,-8.5) -- (34,-8.5);
  \node[lab, anchor=south] at (56,4.2) {\iflangja{要求}{requests}};
  \node[lab, anchor=north] at (56,-11.2) {\iflangja{応答}{responses}};
  \node[lab, anchor=north] at (56,-20.5)
    {\iflangja{チャネル（カーネル内のバッファ）}{channel (kernel buffers)}};
  % the four system calls
  \draw[->, thick] (22,23.5) -- (22,1.5) -- (30,1.5);
  \draw[->, thick] (82,1.5) -- (90,1.5) -- (90,23.5);
  \draw[->, thick] (106,23.5) -- (106,-8.5) -- (82,-8.5);
  \draw[->, thick] (30,-8.5) -- (6,-8.5) -- (6,23.5);
  \node[stepn] at (22,19) {1};
  \node[stepn] at (90,19) {2};
  \node[stepn] at (106,19) {3};
  \node[stepn] at (6,19) {4};
  \node[lab, anchor=west, text=black] at (24.2,13) {\texttt{send}};
  \node[lab, anchor=east, text=black] at (88,13)   {\texttt{recv}};
  \node[lab, anchor=east, text=black] at (104,13)  {\texttt{send}};
  \node[lab, anchor=west, text=black] at (8.2,13)  {\texttt{recv}};
\end{tikzpicture}
   (width ~116 mm)
\begin{tikzpicture}[x=1mm, y=1mm, font=\footnotesize\sffamily,
  proc/.style={draw, rounded corners=2pt, fill=ln-box, align=center,
               minimum width=26mm, minimum height=9mm, inner sep=1pt},
  port/.style={draw, fill=ln-accent!15, minimum width=4mm,
               minimum height=21mm, inner sep=0pt},
  cell/.style={draw, fill=white, minimum width=4mm, minimum height=4.5mm,
               inner sep=0pt},
  msgc/.style={cell, fill=ln-accent!35},
  stepn/.style={circle, draw, fill=white, inner sep=0.5pt,
               minimum size=3.6mm, font=\scriptsize\sffamily\bfseries},
  lab/.style={font=\scriptsize\sffamily, text=ln-muted},
  >={Stealth[length=1.6mm]}]
  % processes
  \node[proc] (p1) at (14,28) {P1 (\iflangja{クライアント}{client})};
  \node[proc] (p2) at (98,28) {P2 (\iflangja{サーバ}{server})};
  % user/kernel boundary
  \draw[dashed, ln-rule, thick] (0,19) -- (116,19);
  \node[lab, anchor=south] at (56,19.5) {\iflangja{ユーザモード}{user mode}};
  \node[lab, anchor=north] at (56,18.5) {\iflangja{カーネルモード}{kernel mode}};
  % channel
  \draw[rounded corners=2pt, fill=ln-box] (26,-20) rectangle (86,10);
  \node[port] (pt1) at (32,-3.5) {};
  \node[port] (pt2) at (80,-3.5) {};
  \node[lab, anchor=north] at (32,-14.6) {\iflangja{ポート}{port}};
  \node[lab, anchor=north] at (80,-14.6) {\iflangja{ポート}{port}};
  \foreach \i in {0,...,7} {
    \pgfmathsetmacro{\x}{42+4*\i}
    \ifnum\i>5 \node[msgc] at (\x,1.5) {}; \else \node[cell] at (\x,1.5) {}; \fi
    \ifnum\i<2 \node[msgc] at (\x,-8.5) {}; \else \node[cell] at (\x,-8.5) {}; \fi
  }
  \draw[->] (34,1.5) -- (40,1.5);
  \draw[->] (72,1.5) -- (78,1.5);
  \draw[->] (78,-8.5) -- (72,-8.5);
  \draw[->] (40,-8.5) -- (34,-8.5);
  \node[lab, anchor=south] at (56,4.2) {\iflangja{要求}{requests}};
  \node[lab, anchor=north] at (56,-11.2) {\iflangja{応答}{responses}};
  \node[lab, anchor=north] at (56,-20.5)
    {\iflangja{チャネル（カーネル内のバッファ）}{channel (kernel buffers)}};
  % the four system calls
  \draw[->, thick] (22,23.5) -- (22,1.5) -- (30,1.5);
  \draw[->, thick] (82,1.5) -- (90,1.5) -- (90,23.5);
  \draw[->, thick] (106,23.5) -- (106,-8.5) -- (82,-8.5);
  \draw[->, thick] (30,-8.5) -- (6,-8.5) -- (6,23.5);
  \node[stepn] at (22,19) {1};
  \node[stepn] at (90,19) {2};
  \node[stepn] at (106,19) {3};
  \node[stepn] at (6,19) {4};
  \node[lab, anchor=west, text=black] at (24.2,13) {\texttt{send}};
  \node[lab, anchor=east, text=black] at (88,13)   {\texttt{recv}};
  \node[lab, anchor=east, text=black] at (104,13)  {\texttt{send}};
  \node[lab, anchor=west, text=black] at (8.2,13)  {\texttt{recv}};
\end{tikzpicture}

  \caption{A request-response exchange over message-based IPC.  Each
    numbered crossing of the user/kernel boundary is one system call and one
    copy between a process's address space and the kernel's channel.}
  \label{fig:l09-message-ipc}
\end{figure}

The advantage of message passing is its simplicity: the kernel performs the
channel management and the synchronization---a receive on an empty channel,
for instance, blocks until a message arrives.  The disadvantage is overhead:
the crossings and copies are paid for every message.

\section{Forms of message passing}
\label{sec:l09-forms}

Unix-like systems provide three common forms of message
passing,\source{P3L3, forms of message passing; OSTEP ch.~5; Kerrisk ch.~44, 46,
56.}
which differ in what they transfer and in how much processing the kernel
performs.

\begin{definition}[Pipe]
  A \term{pipe}[パイプ] is a channel between two endpoints that carries a
  byte stream: bytes written into one end are read from the other end in
  the same order, without message boundaries.
\end{definition}

The POSIX call \texttt{pipe(fd)} creates a pipe and stores one file
descriptor for each end in the two-element array \texttt{fd}.  A typical use connects the output of one process
to the input of another, as the shell does for \texttt{ls | wc -l}.

\begin{definition}[Message queue]
  A \term{message queue}[メッセージキュー] is a channel that carries
  discrete messages among processes: a sender submits a properly formatted
  message, and a receiver obtains complete messages.
\end{definition}

Because the kernel knows where messages begin and end, it can manage them,
for example assign priorities to messages and schedule their delivery
accordingly.  Unix-like systems offer a System~V and a POSIX interface for
message queues.

\begin{definition}[Socket]
  A \term{socket}[ソケット] is a communication endpoint created with the
  \texttt{socket()} system call, which allocates a kernel-level socket
  buffer and associates with it the kernel-level processing that the
  communication needs, such as a TCP/IP protocol stack.
\end{definition}

The calls \texttt{send()} and \texttt{recv()} pass message buffers into the
socket buffer and out of it.  If the processes run on different machines,
the channel connects the process, through the protocol stack, to a network
device; if they run on the same machine, the kernel can bypass the full
protocol stack.  \Cref{tab:l09-forms} compares the three forms.

\begin{table}[htbp]
  \centering
  \caption{The three forms of message passing.}
  \label{tab:l09-forms}
  \small
  \begin{tabular}{@{}llll@{}}
    \toprule
    & Pipe & Message queue & Socket \\
    \midrule
    Unit of transfer & byte stream & message & message buffer \\
    Kernel processing & buffering & priorities, delivery & protocol stack \\
    Machines & same & same & same or different \\
    Interface & \texttt{pipe()} & System~V, POSIX & \texttt{socket()} \\
    \bottomrule
  \end{tabular}
\end{table}

\section{Shared-memory IPC}
\label{sec:l09-shm}

In shared-memory IPC the processes read and write a region of memory that
all of them can access.\source{P3L3, shared memory IPC; Silberschatz et al.\
ch.~3.}\index{shared memory}  The operating system establishes this channel
by mapping the same physical pages into the virtual address spaces of the
processes: page-table entries of every participating process point to the
same page frames (\cref{fig:l09-shm-mapping}).  As with any paged mapping
(Lesson~8), the virtual addresses at which the processes see the region
need not be the same, and the physical frames need not be contiguous.

\begin{figure}[htbp]
  \centering
  % Shared-memory IPC: the same (non-contiguous) physical frames are mapped
% into two virtual address spaces at different virtual addresses.
% Usage: % Shared-memory IPC: the same (non-contiguous) physical frames are mapped
% into two virtual address spaces at different virtual addresses.
% Usage: % Shared-memory IPC: the same (non-contiguous) physical frames are mapped
% into two virtual address spaces at different virtual addresses.
% Usage: \input{figures/l09-shm-mapping}   (width ~116 mm)
\begin{tikzpicture}[x=1mm, y=1mm, font=\footnotesize\sffamily,
  slot/.style={draw, fill=white, minimum width=18mm, minimum height=5mm,
               inner sep=0pt, font=\scriptsize\sffamily},
  sh/.style={slot, fill=ln-accent!20},
  oth/.style={slot, fill=ln-box},
  hd/.style={font=\scriptsize\sffamily, align=center, text width=30mm},
  lab/.style={font=\scriptsize\sffamily, text=ln-muted},
  >={Stealth[length=1.6mm]}]
  \foreach \i in {0,...,7} {
    \node[slot] (a\i) at (17, 2.5+5*\i) {};
    \node[slot] (m\i) at (59, 2.5+5*\i) {};
    \node[slot] (b\i) at (101, 2.5+5*\i) {};
  }
  % other (private) frames
  \foreach \i in {0,3,4,7} { \node[oth] at (59, 2.5+5*\i) {}; }
  % shared pages
  \node[sh] (a5) at (17,27.5)  {\iflangja{共有 0}{shared 0}};
  \node[sh] (a6) at (17,32.5)  {\iflangja{共有 1}{shared 1}};
  \node[sh] (m6) at (59,32.5)  {\iflangja{共有 0}{shared 0}};
  \node[sh] (m2) at (59,12.5)  {\iflangja{共有 1}{shared 1}};
  \node[sh] (b1) at (101,7.5)  {\iflangja{共有 0}{shared 0}};
  \node[sh] (b2) at (101,12.5) {\iflangja{共有 1}{shared 1}};
  \draw[->] (a5.east) -- (m6.west);
  \draw[->] (a6.east) -- (m2.west);
  \draw[->] (b1.west) -- (m6.east);
  \draw[->] (b2.west) -- (m2.east);
  % headers
  \node[hd, anchor=south] at (17,41)  {\iflangja{P1 の仮想アドレス空間}{virtual address\\space of P1}};
  \node[hd, anchor=south] at (59,41)  {\iflangja{物理メモリ}{physical memory}};
  \node[hd, anchor=south] at (101,41) {\iflangja{P2 の仮想アドレス空間}{virtual address\\space of P2}};
  % virtual addresses of the region
  \draw[ln-muted] (6,25) -- (8,25);
  \node[lab, anchor=east] at (6,25) {$\mathit{VA}_1$};
  \draw[ln-muted] (110,5) -- (112,5);
  \node[lab, anchor=west] at (112,5) {$\mathit{VA}_2$};
\end{tikzpicture}
   (width ~116 mm)
\begin{tikzpicture}[x=1mm, y=1mm, font=\footnotesize\sffamily,
  slot/.style={draw, fill=white, minimum width=18mm, minimum height=5mm,
               inner sep=0pt, font=\scriptsize\sffamily},
  sh/.style={slot, fill=ln-accent!20},
  oth/.style={slot, fill=ln-box},
  hd/.style={font=\scriptsize\sffamily, align=center, text width=30mm},
  lab/.style={font=\scriptsize\sffamily, text=ln-muted},
  >={Stealth[length=1.6mm]}]
  \foreach \i in {0,...,7} {
    \node[slot] (a\i) at (17, 2.5+5*\i) {};
    \node[slot] (m\i) at (59, 2.5+5*\i) {};
    \node[slot] (b\i) at (101, 2.5+5*\i) {};
  }
  % other (private) frames
  \foreach \i in {0,3,4,7} { \node[oth] at (59, 2.5+5*\i) {}; }
  % shared pages
  \node[sh] (a5) at (17,27.5)  {\iflangja{共有 0}{shared 0}};
  \node[sh] (a6) at (17,32.5)  {\iflangja{共有 1}{shared 1}};
  \node[sh] (m6) at (59,32.5)  {\iflangja{共有 0}{shared 0}};
  \node[sh] (m2) at (59,12.5)  {\iflangja{共有 1}{shared 1}};
  \node[sh] (b1) at (101,7.5)  {\iflangja{共有 0}{shared 0}};
  \node[sh] (b2) at (101,12.5) {\iflangja{共有 1}{shared 1}};
  \draw[->] (a5.east) -- (m6.west);
  \draw[->] (a6.east) -- (m2.west);
  \draw[->] (b1.west) -- (m6.east);
  \draw[->] (b2.west) -- (m2.east);
  % headers
  \node[hd, anchor=south] at (17,41)  {\iflangja{P1 の仮想アドレス空間}{virtual address\\space of P1}};
  \node[hd, anchor=south] at (59,41)  {\iflangja{物理メモリ}{physical memory}};
  \node[hd, anchor=south] at (101,41) {\iflangja{P2 の仮想アドレス空間}{virtual address\\space of P2}};
  % virtual addresses of the region
  \draw[ln-muted] (6,25) -- (8,25);
  \node[lab, anchor=east] at (6,25) {$\mathit{VA}_1$};
  \draw[ln-muted] (110,5) -- (112,5);
  \node[lab, anchor=west] at (112,5) {$\mathit{VA}_2$};
\end{tikzpicture}
   (width ~116 mm)
\begin{tikzpicture}[x=1mm, y=1mm, font=\footnotesize\sffamily,
  slot/.style={draw, fill=white, minimum width=18mm, minimum height=5mm,
               inner sep=0pt, font=\scriptsize\sffamily},
  sh/.style={slot, fill=ln-accent!20},
  oth/.style={slot, fill=ln-box},
  hd/.style={font=\scriptsize\sffamily, align=center, text width=30mm},
  lab/.style={font=\scriptsize\sffamily, text=ln-muted},
  >={Stealth[length=1.6mm]}]
  \foreach \i in {0,...,7} {
    \node[slot] (a\i) at (17, 2.5+5*\i) {};
    \node[slot] (m\i) at (59, 2.5+5*\i) {};
    \node[slot] (b\i) at (101, 2.5+5*\i) {};
  }
  % other (private) frames
  \foreach \i in {0,3,4,7} { \node[oth] at (59, 2.5+5*\i) {}; }
  % shared pages
  \node[sh] (a5) at (17,27.5)  {\iflangja{共有 0}{shared 0}};
  \node[sh] (a6) at (17,32.5)  {\iflangja{共有 1}{shared 1}};
  \node[sh] (m6) at (59,32.5)  {\iflangja{共有 0}{shared 0}};
  \node[sh] (m2) at (59,12.5)  {\iflangja{共有 1}{shared 1}};
  \node[sh] (b1) at (101,7.5)  {\iflangja{共有 0}{shared 0}};
  \node[sh] (b2) at (101,12.5) {\iflangja{共有 1}{shared 1}};
  \draw[->] (a5.east) -- (m6.west);
  \draw[->] (a6.east) -- (m2.west);
  \draw[->] (b1.west) -- (m6.east);
  \draw[->] (b2.west) -- (m2.east);
  % headers
  \node[hd, anchor=south] at (17,41)  {\iflangja{P1 の仮想アドレス空間}{virtual address\\space of P1}};
  \node[hd, anchor=south] at (59,41)  {\iflangja{物理メモリ}{physical memory}};
  \node[hd, anchor=south] at (101,41) {\iflangja{P2 の仮想アドレス空間}{virtual address\\space of P2}};
  % virtual addresses of the region
  \draw[ln-muted] (6,25) -- (8,25);
  \node[lab, anchor=east] at (6,25) {$\mathit{VA}_1$};
  \draw[ln-muted] (110,5) -- (112,5);
  \node[lab, anchor=west] at (112,5) {$\mathit{VA}_2$};
\end{tikzpicture}

  \caption{Shared memory.  Two page frames are mapped into both address
    spaces; P1 sees the region at virtual address $\mathit{VA}_1$, P2 at a
    different address $\mathit{VA}_2$, and the frames are not adjacent in
    physical memory.}
  \label{fig:l09-shm-mapping}
\end{figure}

\caution{Since in general $\mathit{VA}_1 \neq \mathit{VA}_2$, a pointer that
one process stores in shared memory is meaningless to the other; shared data
structures should use offsets from the start of the region.}

Once the mapping exists, the operating system is out of the way: reads and
writes are ordinary memory accesses, and system calls are needed only for
the setup.  Data copies are potentially reduced, but not eliminated.  Data
that a process creates in the shared region are immediately visible to the
others, whereas data that already exist elsewhere in its address space must
still be copied into the region.

The price is that the operating system no longer performs the services it
provides for message passing.  The processes must synchronize their
concurrent accesses explicitly (\cref{sec:l09-sync}), and the communication
protocol---the format of the data, how a receiver learns that data are
available, how the space in the region is managed---is the programmer's
responsibility (\cref{sec:l09-design}).  Unix and Linux support shared
memory through System~V shared memory (\cref{sec:l09-sysv}), POSIX shared
memory (\cref{sec:l09-posix}) and memory-mapped files; Android adds its own
facility, ashmem.

\section{Copy versus map}
\label{sec:l09-copy-map}

Message passing and shared memory pursue the same goal, transferring data
from one address space into another, but spend CPU time on different
operations.\source{P3L3, copy vs.\ map; Silberschatz et al.\ ch.~3.}
\begin{description}
  \item[Copy (message passing)] CPU time is spent on copying the data into
    the port and out of it again, with a system call for every send and
    receive---for every message.
  \item[Map (shared memory)] CPU time is spent on mapping physical memory
    into the address spaces---once, when the region is set up---and on
    copying data into the region whenever they are not created there.
\end{description}

Mapping is relatively expensive.  Set up once and used for many transfers,
it pays off, as the example in Lesson~2 showed.  For large data it can pay
off even for a single transfer, because copying a large amount of data from
one address space into another takes far longer than mapping it:
$t_{\mathrm{copy}} \gg t_{\mathrm{map}}$.  A simple cost model makes both
statements precise.  Let a system call cost $t_{\mathrm{sys}}$, copying cost
$c$ per byte, and let $n$ transfers of $S$ bytes each go from one process to
another.  Message passing needs two system calls and two copies per
transfer.  Shared memory needs a one-time setup of cost $t_{\mathrm{setup}}$
and one copy per transfer, into the region, while the receiver reads the
data in place:
\begin{equation}
  T_{\mathrm{copy}}(n) = n\,(2\,t_{\mathrm{sys}} + 2cS) ,
  \qquad
  T_{\mathrm{map}}(n) = t_{\mathrm{setup}} + n\,cS .
  \label{eq:l09-copy-map}
\end{equation}
Shared memory is cheaper for $n > n^{*} = t_{\mathrm{setup}} /
(2\,t_{\mathrm{sys}} + cS)$.  The setup cost consists of a fixed part and a
part proportional to the number of pages mapped; as long as mapping a page
costs less than copying its contents, $n^{*}$ falls as $S$ grows and drops
below one for large data.

\begin{example}
  \label{ex:l09-copy-map}
  Let $t_{\mathrm{sys}} = \SI{1}{\micro\second}$,
  $c = \SI{0.1}{\micro\second}$ per kilobyte, and let the setup cost
  \SI{50}{\micro\second} plus \SI{0.1}{\micro\second} per \SI{4}{\kilo\byte}
  page.
  \begin{itemize}
    \item $S = \SI{4}{\kilo\byte}$ (one page): a transfer costs
      $2 + 2 \times 0.4 = \SI{2.8}{\micro\second}$ by message passing and
      \SI{0.4}{\micro\second} over shared memory after a setup of
      \SI{50.1}{\micro\second}, so $n^{*} = 50.1/2.4 \approx 20.9$.  From
      21 transfers on, shared memory is cheaper.
    \item $S = \SI{40}{\mega\byte}$ ($10^4$ pages): a single transfer costs
      $2 + 2 \times 4000 = \SI{8002}{\micro\second}$ by message passing, but
      $1050 + 4000 = \SI{5050}{\micro\second}$ over shared memory, and only
      the \SI{1050}{\micro\second} of setup if the sender creates the data
      directly in the region; $n^{*} = 1050/4002 \approx 0.26$.
  \end{itemize}
\end{example}

Windows exploits this trade-off in its facility for communication between
processes on the same machine, the
\term{local procedure call}[ローカル手続き呼び出し]
(LPC).\index{LPC|see{local procedure call}}  Small messages are copied
through a port; for larger data, a shared memory region is mapped into both
processes, and the data are exchanged through it.

\begin{keypoint}
  Message passing is simple---the kernel manages the channel and
  synchronizes the processes---but every message pays for system calls and
  two copies.  Shared memory pays a one-time mapping cost and then bypasses
  the kernel, at the price of explicit synchronization and a
  programmer-defined protocol.  Copy small messages; map large data or
  regions reused for many transfers.
\end{keypoint}

\section{System V shared memory}
\label{sec:l09-sysv}

The System~V interface,\source{P3L3, SysV shared memory, SysV shared
memory API; Kerrisk ch.~45, 48.} one of the oldest Unix IPC interfaces,
organizes shared memory in segments.

\begin{definition}[Shared memory segment]
  A \term{shared memory segment}[共有メモリセグメント] is a region of
  physical memory, not necessarily made of contiguous pages, that the
  operating system manages as a system-wide object and that any number of
  processes can map into their address spaces.
\end{definition}

Because segments are system-wide, the operating system limits both their
number and their total size.\margin{On Linux the limits are the kernel
parameters \texttt{shmmni}, \texttt{shmmax} and \texttt{shmall}; by default
at most 4096 segments may exist.}  A segment is identified by a
\emph{key}: every process that uses the same key accesses the same segment.
Its life cycle consists of four operations (\cref{fig:l09-sysv-lifecycle}).
\begin{description}
  \item[Create] The operating system allocates the memory of the segment and
    associates it with the key.
  \item[Attach] The operating system maps the segment into the address space
    of the calling process, i.e., it establishes valid mappings from virtual
    addresses of the process to the physical memory of the segment.  Several
    processes may attach the same segment, each at its own virtual address.
  \item[Detach] The operating system invalidates these mappings: the
    page-table entries for the virtual address range of the segment are no
    longer valid.  The segment itself remains.
  \item[Destroy] Only an explicit request removes a segment.  Until then it
    persists, even when no process has it attached, and it may be attached
    and detached many times.
\end{description}

\begin{figure}[htbp]
  \centering
  % Life cycle of a System V shared memory segment.
% Usage: % Life cycle of a System V shared memory segment.
% Usage: % Life cycle of a System V shared memory segment.
% Usage: \input{figures/l09-sysv-lifecycle}   (width ~120 mm)
\begin{tikzpicture}[x=1mm, y=1mm, font=\footnotesize\sffamily,
  st/.style={draw, rounded corners=3pt, fill=ln-box, align=center,
             minimum height=11mm, inner sep=2pt},
  api/.style={font=\scriptsize\ttfamily, align=center, inner sep=1pt},
  >={Stealth[length=1.8mm]}]
  \node[st, text width=17mm] (none) at (0,0)
    {\iflangja{セグメント\\なし}{no\\segment}};
  \node[st, text width=24mm] (free) at (48,0)
    {\iflangja{存在する\\（接続なし）}{exists,\\not attached}};
  \node[st, text width=24mm] (att) at (96,0)
    {\iflangja{$n\ge 1$ 個のプロセス\\が接続}{attached by\\$n\ge 1$ processes}};
  \node[st, text width=24mm] (mark) at (96,-30)
    {\iflangja{破棄予約済み}{marked for\\destruction}};
  \draw[->] (none) -- node[api, above] {shmget()}
                      node[api, below] {IPC\_CREAT} (free);
  \draw[->] ([yshift=2.5mm]free.east) -- node[api, above] {shmat()}
                                         ([yshift=2.5mm]att.west);
  \draw[->] ([yshift=-2.5mm]att.west) -- node[api, below] {shmdt()\\(\iflangja{最後}{last})}
                                         ([yshift=-2.5mm]free.east);
  \draw[->] ([xshift=-4mm]att.north) .. controls +(-3,12) and +(3,12) ..
    node[api, above] {shmat() / shmdt()} ([xshift=4mm]att.north);
  \draw[->] (free.north) .. controls +(0,12) and +(0,12) ..
    node[api, above] {shmctl(IPC\_RMID)} (none.north);
  \draw[->] (att) -- node[api, left=1mm] {shmctl(IPC\_RMID)} (mark);
  \draw[->] (mark.west) -| node[api, above, pos=0.3] {shmdt()\,(\iflangja{最後}{last})} (none.south);
\end{tikzpicture}
   (width ~120 mm)
\begin{tikzpicture}[x=1mm, y=1mm, font=\footnotesize\sffamily,
  st/.style={draw, rounded corners=3pt, fill=ln-box, align=center,
             minimum height=11mm, inner sep=2pt},
  api/.style={font=\scriptsize\ttfamily, align=center, inner sep=1pt},
  >={Stealth[length=1.8mm]}]
  \node[st, text width=17mm] (none) at (0,0)
    {\iflangja{セグメント\\なし}{no\\segment}};
  \node[st, text width=24mm] (free) at (48,0)
    {\iflangja{存在する\\（接続なし）}{exists,\\not attached}};
  \node[st, text width=24mm] (att) at (96,0)
    {\iflangja{$n\ge 1$ 個のプロセス\\が接続}{attached by\\$n\ge 1$ processes}};
  \node[st, text width=24mm] (mark) at (96,-30)
    {\iflangja{破棄予約済み}{marked for\\destruction}};
  \draw[->] (none) -- node[api, above] {shmget()}
                      node[api, below] {IPC\_CREAT} (free);
  \draw[->] ([yshift=2.5mm]free.east) -- node[api, above] {shmat()}
                                         ([yshift=2.5mm]att.west);
  \draw[->] ([yshift=-2.5mm]att.west) -- node[api, below] {shmdt()\\(\iflangja{最後}{last})}
                                         ([yshift=-2.5mm]free.east);
  \draw[->] ([xshift=-4mm]att.north) .. controls +(-3,12) and +(3,12) ..
    node[api, above] {shmat() / shmdt()} ([xshift=4mm]att.north);
  \draw[->] (free.north) .. controls +(0,12) and +(0,12) ..
    node[api, above] {shmctl(IPC\_RMID)} (none.north);
  \draw[->] (att) -- node[api, left=1mm] {shmctl(IPC\_RMID)} (mark);
  \draw[->] (mark.west) -| node[api, above, pos=0.3] {shmdt()\,(\iflangja{最後}{last})} (none.south);
\end{tikzpicture}
   (width ~120 mm)
\begin{tikzpicture}[x=1mm, y=1mm, font=\footnotesize\sffamily,
  st/.style={draw, rounded corners=3pt, fill=ln-box, align=center,
             minimum height=11mm, inner sep=2pt},
  api/.style={font=\scriptsize\ttfamily, align=center, inner sep=1pt},
  >={Stealth[length=1.8mm]}]
  \node[st, text width=17mm] (none) at (0,0)
    {\iflangja{セグメント\\なし}{no\\segment}};
  \node[st, text width=24mm] (free) at (48,0)
    {\iflangja{存在する\\（接続なし）}{exists,\\not attached}};
  \node[st, text width=24mm] (att) at (96,0)
    {\iflangja{$n\ge 1$ 個のプロセス\\が接続}{attached by\\$n\ge 1$ processes}};
  \node[st, text width=24mm] (mark) at (96,-30)
    {\iflangja{破棄予約済み}{marked for\\destruction}};
  \draw[->] (none) -- node[api, above] {shmget()}
                      node[api, below] {IPC\_CREAT} (free);
  \draw[->] ([yshift=2.5mm]free.east) -- node[api, above] {shmat()}
                                         ([yshift=2.5mm]att.west);
  \draw[->] ([yshift=-2.5mm]att.west) -- node[api, below] {shmdt()\\(\iflangja{最後}{last})}
                                         ([yshift=-2.5mm]free.east);
  \draw[->] ([xshift=-4mm]att.north) .. controls +(-3,12) and +(3,12) ..
    node[api, above] {shmat() / shmdt()} ([xshift=4mm]att.north);
  \draw[->] (free.north) .. controls +(0,12) and +(0,12) ..
    node[api, above] {shmctl(IPC\_RMID)} (none.north);
  \draw[->] (att) -- node[api, left=1mm] {shmctl(IPC\_RMID)} (mark);
  \draw[->] (mark.west) -| node[api, above, pos=0.3] {shmdt()\,(\iflangja{最後}{last})} (none.south);
\end{tikzpicture}

  \caption{Life cycle of a System~V shared memory segment.  A segment
    outlives the processes that use it and is removed only by
    \texttt{IPC\_RMID}; a segment removed while still attached is destroyed
    when the last process detaches.}
  \label{fig:l09-sysv-lifecycle}
\end{figure}

A segment therefore behaves like a file rather than like the private memory
of a process: a process that exits is detached from its segments, but the
segments remain until they are destroyed or the system
reboots.\caution{A program that exits without destroying its segments
leaks their memory until reboot; \cref{sec:l09-tools} shows how to remove
such segments.}

\subsection{The System V API}

The interface consists of four calls and a helper for generating keys.
\begin{description}
  \item[\texttt{shmget(key, size, flags)}] creates a segment of the given
    size if \texttt{flags} contains \texttt{IPC\_CREAT}, or opens the
    existing segment with this key, and returns its identifier.
  \item[\texttt{ftok(pathname, proj\_id)}] generates a key from the name of
    an existing file and a small integer.  The same arguments always yield
    the same key, so unrelated processes can find a segment by agreeing on a
    file name and a number.
  \item[\texttt{shmat(id, addr, flags)}] attaches the segment and returns the
    virtual address at which it is mapped; with \texttt{addr} set to
    \texttt{NULL} the operating system chooses the address.  The returned
    untyped pointer can be cast to any data structure placed in the segment.
  \item[\texttt{shmdt(addr)}] detaches the segment mapped at \texttt{addr}.
  \item[\texttt{shmctl(id, cmd, buf)}] passes a control command for the
    segment to the operating system; \texttt{IPC\_RMID} destroys it.
\end{description}

\begin{example}
  \label{ex:l09-sysv}
  Process A creates a segment and writes a string into it; process B, which
  runs later, finds the segment through the same key, prints the string and
  destroys the segment (error checks omitted).
  \begin{lstlisting}[language=C, numbers=none]
/* process A */
key_t key = ftok("/tmp/app", 'A');
int id = shmget(key, 4096, IPC_CREAT | 0600);
char *p = shmat(id, NULL, 0);
strcpy(p, "hello");                /* visible to B  */
shmdt(p);                          /* segment stays */

/* process B */
key_t key = ftok("/tmp/app", 'A'); /* same key      */
int id = shmget(key, 4096, 0);     /* open          */
char *q = shmat(id, NULL, 0);
printf("%s\n", q);
shmdt(q);
shmctl(id, IPC_RMID, NULL);        /* destroy       */
  \end{lstlisting}
  Between A's \texttt{shmdt} and B's \texttt{shmat} no process has the
  segment attached, yet its contents are preserved.
\end{example}

\section{POSIX shared memory}
\label{sec:l09-posix}

POSIX standardizes a second shared memory interface, which is not as widely
supported as System~V; Linux, for instance, has supported it since kernel
2.4.\source{P3L3, POSIX shared memory API; Kerrisk ch.~49, 54.}  Instead
of segments identified by keys it uses files, identified by name, that exist
only in a memory-backed temporary file system (tmpfs), and the processes
access them through the general mechanism for mapping files into
memory.\margin{On Linux these files appear in \texttt{/dev/shm}.}

\begin{definition}[Memory-mapped file]
  A \term{memory-mapped file}[メモリマップドファイル] is a file whose
  contents are mapped into the virtual address space of a process, so that
  the process accesses the file by reading and writing memory.  Processes
  that map the same file as shared map the same physical pages and see each
  other's changes.
\end{definition}

A memory-mapped file on an ordinary file system is itself an IPC mechanism
whose data also persist in the file; POSIX shared memory applies the same
mechanism to a file that exists only in memory.  Its operations correspond
to those of System~V (\cref{tab:l09-shm-api}).
\texttt{shm\_open(name, flags, mode)} creates or opens the shared memory
object and returns a file descriptor for it, and \texttt{ftruncate} sets
its size.  \texttt{mmap} with the flag \texttt{MAP\_SHARED} attaches the
object and returns its virtual address; \texttt{munmap} detaches it.
\texttt{close} releases the file descriptor without affecting an
established mapping, and \texttt{shm\_unlink(name)} removes the name; the
object is destroyed once no process has it open or mapped.

\begin{table}[htbp]
  \centering
  \caption{Corresponding operations of the two shared memory interfaces.}
  \label{tab:l09-shm-api}
  \small
  \begin{tabular}{@{}lll@{}}
    \toprule
    Operation & System~V & POSIX \\
    \midrule
    Name the region & key from \texttt{ftok()} & file name in tmpfs \\
    Create or open & \texttt{shmget()} & \texttt{shm\_open()}, \texttt{ftruncate()} \\
    Attach & \texttt{shmat()} & \texttt{mmap()} \\
    Detach & \texttt{shmdt()} & \texttt{munmap()} \\
    Destroy & \texttt{shmctl(IPC\_RMID)} & \texttt{shm\_unlink()} \\
    \bottomrule
  \end{tabular}
\end{table}

\section{Shared memory and synchronization}
\label{sec:l09-sync}

Processes that access a shared region concurrently are in the situation of
the threads of Lesson~3, which share state within one address space: without
synchronization their accesses form race conditions.\source{P3L3, shared
memory and synchronization; Kerrisk ch.~45--47.}  Two kinds of mechanisms
are available: those of a threading library, such as the Pthreads mutexes
and condition variables, configured for use across processes; and
synchronization based on IPC mechanisms that the operating system provides,
such as message queues and semaphores.  With either kind the processes must
coordinate two things: the number of concurrent accesses to the shared
region, which mutual exclusion limits, and the moment at which data become
available for consumption, which one process must be able to wait for and
the other to announce, as with a condition variable.

\subsection{Pthreads synchronization across processes}

The Pthreads mutexes and condition variables of Lesson~4 can synchronize
processes under two conditions.\source{P3L3, PThreads sync for IPC.}
First, the attribute objects with which they are initialized,
\texttt{pthread\_mutexattr\_t} and \texttt{pthread\_condattr\_t}, must have
the process-shared attribute set to \texttt{PTHREAD\_PROCESS\_SHARED}; by
default only the threads of the creating process may use them.  Second, the
mutex and the condition variable themselves must reside in the shared
region, so that all processes operate on the same data structures.

\needspace{9\baselineskip}
\begin{example}
  \label{ex:l09-pshared}
  A process defines the layout of the region, which begins with the
  synchronization constructs:
  \begin{lstlisting}[language=C, numbers=none]
typedef struct {
  pthread_mutex_t mutex;
  pthread_cond_t  cond;
  char            data[1024];
} shm_data_t;
  \end{lstlisting}
  \needspace{6\baselineskip}
  It creates and attaches the segment:
  \begin{lstlisting}[language=C, numbers=none]
key_t key = ftok("/tmp/app", 120);
int id = shmget(key, sizeof(shm_data_t),
                IPC_CREAT | IPC_EXCL | 0600);
shm_data_t *shm = shmat(id, NULL, 0);
  \end{lstlisting}
  \needspace{7\baselineskip}
  and it initializes a process-shared mutex inside the region:
  \begin{lstlisting}[language=C, numbers=none]
pthread_mutexattr_t attr;
pthread_mutexattr_init(&attr);
pthread_mutexattr_setpshared(&attr,
                             PTHREAD_PROCESS_SHARED);
pthread_mutex_init(&shm->mutex, &attr);
  \end{lstlisting}
  The condition variable is initialized likewise, with the
  \texttt{pthread\_condattr} counterparts of these calls.  The flag
  \texttt{IPC\_EXCL} makes \texttt{shmget} fail if the segment already
  exists, so the synchronization constructs are initialized only once.
  Every process that attaches the segment then locks \texttt{shm->mutex}
  and waits on \texttt{shm->cond} exactly as threads would.
\end{example}

\subsection{Synchronization with other IPC mechanisms}

Processes can also synchronize through the IPC facilities of the operating
system.\source{P3L3, sync for other IPC; Kerrisk ch.~46--47.}  A message
queue can provide mutual exclusion through a protocol of sends and receives
(\cref{fig:l09-mq-sync}).  P1 writes data into the shared region and then
sends a \enquote{ready} message;\margin{When both directions use one System~V
queue, each process selects its peer's messages by the message type
argument of \texttt{msgrcv()}.} P2, blocked in a receive on the queue,
reads the data when that message arrives and replies with an \enquote{ok}
message, which P1 awaits before it accesses the region again.  Because each
process accesses the region only after receiving the other's message, the
accesses never overlap.

\begin{figure}[htbp]
  \centering
  % Using a message queue to coordinate access to shared memory.
% Usage: % Using a message queue to coordinate access to shared memory.
% Usage: % Using a message queue to coordinate access to shared memory.
% Usage: \input{figures/l09-mq-sync}   (width ~116 mm)
\begin{tikzpicture}[x=1mm, y=1mm, font=\footnotesize\sffamily,
  hd/.style={draw, rounded corners=2pt, fill=ln-box, align=center,
             minimum height=9mm, inner sep=1.5pt,
             text height=1.6ex, text depth=0.5ex},
  msg/.style={font=\scriptsize\sffamily, fill=white, inner sep=0.6pt, above=0.5mm},
  lab/.style={font=\scriptsize\sffamily, text=ln-muted},
  bar/.style={draw=ln-accent, fill=ln-accent!20},
  >={Stealth[length=1.6mm]}]
  \node[hd, minimum width=16mm] (h1) at (10,0) {P1};
  \node[hd, minimum width=25mm] (hs) at (44,0) {\iflangja{共有メモリ}{shared memory}};
  \node[hd, minimum width=25mm] (hq) at (78,0) {\iflangja{メッセージキュー}{message queue}};
  \node[hd, minimum width=16mm] (h2) at (108,0) {P2};
  \foreach \x/\n in {10/h1,44/hs,78/hq,108/h2}
    \draw[dashed, ln-rule] (\n.south) -- (\x,-62);
  % blocking periods
  \draw[bar] (107,-6) rectangle (109,-28);
  \node[lab, anchor=east] at (105.5,-7) {\iflangja{\texttt{msgrcv()} で待つ}{blocked in \texttt{msgrcv()}}};
  \draw[bar] (9,-21) rectangle (11,-55);
  \node[lab, anchor=west] at (12.5,-24) {\iflangja{\texttt{msgrcv()} で待つ}{blocked in \texttt{msgrcv()}}};
  % messages
  \draw[->] (10,-12) -- node[msg] {1\ \iflangja{データを書く}{write data}} (44,-12);
  \draw[->] (10,-19) -- node[msg, pos=0.62] {2\ \texttt{msgsnd(ready)}} (78,-19);
  \draw[->] (78,-28) -- node[msg] {3\ \texttt{ready}} (108,-28);
  \draw[->] (108,-37) -- node[msg, pos=0.6] {4\ \iflangja{データを読む}{read data}} (44,-37);
  \draw[->] (108,-46) -- node[msg] {5\ \texttt{msgsnd(ok)}} (78,-46);
  \draw[->] (78,-55) -- node[msg, pos=0.45] {6\ \texttt{ok}} (10,-55);
\end{tikzpicture}
   (width ~116 mm)
\begin{tikzpicture}[x=1mm, y=1mm, font=\footnotesize\sffamily,
  hd/.style={draw, rounded corners=2pt, fill=ln-box, align=center,
             minimum height=9mm, inner sep=1.5pt,
             text height=1.6ex, text depth=0.5ex},
  msg/.style={font=\scriptsize\sffamily, fill=white, inner sep=0.6pt, above=0.5mm},
  lab/.style={font=\scriptsize\sffamily, text=ln-muted},
  bar/.style={draw=ln-accent, fill=ln-accent!20},
  >={Stealth[length=1.6mm]}]
  \node[hd, minimum width=16mm] (h1) at (10,0) {P1};
  \node[hd, minimum width=25mm] (hs) at (44,0) {\iflangja{共有メモリ}{shared memory}};
  \node[hd, minimum width=25mm] (hq) at (78,0) {\iflangja{メッセージキュー}{message queue}};
  \node[hd, minimum width=16mm] (h2) at (108,0) {P2};
  \foreach \x/\n in {10/h1,44/hs,78/hq,108/h2}
    \draw[dashed, ln-rule] (\n.south) -- (\x,-62);
  % blocking periods
  \draw[bar] (107,-6) rectangle (109,-28);
  \node[lab, anchor=east] at (105.5,-7) {\iflangja{\texttt{msgrcv()} で待つ}{blocked in \texttt{msgrcv()}}};
  \draw[bar] (9,-21) rectangle (11,-55);
  \node[lab, anchor=west] at (12.5,-24) {\iflangja{\texttt{msgrcv()} で待つ}{blocked in \texttt{msgrcv()}}};
  % messages
  \draw[->] (10,-12) -- node[msg] {1\ \iflangja{データを書く}{write data}} (44,-12);
  \draw[->] (10,-19) -- node[msg, pos=0.62] {2\ \texttt{msgsnd(ready)}} (78,-19);
  \draw[->] (78,-28) -- node[msg] {3\ \texttt{ready}} (108,-28);
  \draw[->] (108,-37) -- node[msg, pos=0.6] {4\ \iflangja{データを読む}{read data}} (44,-37);
  \draw[->] (108,-46) -- node[msg] {5\ \texttt{msgsnd(ok)}} (78,-46);
  \draw[->] (78,-55) -- node[msg, pos=0.45] {6\ \texttt{ok}} (10,-55);
\end{tikzpicture}
   (width ~116 mm)
\begin{tikzpicture}[x=1mm, y=1mm, font=\footnotesize\sffamily,
  hd/.style={draw, rounded corners=2pt, fill=ln-box, align=center,
             minimum height=9mm, inner sep=1.5pt,
             text height=1.6ex, text depth=0.5ex},
  msg/.style={font=\scriptsize\sffamily, fill=white, inner sep=0.6pt, above=0.5mm},
  lab/.style={font=\scriptsize\sffamily, text=ln-muted},
  bar/.style={draw=ln-accent, fill=ln-accent!20},
  >={Stealth[length=1.6mm]}]
  \node[hd, minimum width=16mm] (h1) at (10,0) {P1};
  \node[hd, minimum width=25mm] (hs) at (44,0) {\iflangja{共有メモリ}{shared memory}};
  \node[hd, minimum width=25mm] (hq) at (78,0) {\iflangja{メッセージキュー}{message queue}};
  \node[hd, minimum width=16mm] (h2) at (108,0) {P2};
  \foreach \x/\n in {10/h1,44/hs,78/hq,108/h2}
    \draw[dashed, ln-rule] (\n.south) -- (\x,-62);
  % blocking periods
  \draw[bar] (107,-6) rectangle (109,-28);
  \node[lab, anchor=east] at (105.5,-7) {\iflangja{\texttt{msgrcv()} で待つ}{blocked in \texttt{msgrcv()}}};
  \draw[bar] (9,-21) rectangle (11,-55);
  \node[lab, anchor=west] at (12.5,-24) {\iflangja{\texttt{msgrcv()} で待つ}{blocked in \texttt{msgrcv()}}};
  % messages
  \draw[->] (10,-12) -- node[msg] {1\ \iflangja{データを書く}{write data}} (44,-12);
  \draw[->] (10,-19) -- node[msg, pos=0.62] {2\ \texttt{msgsnd(ready)}} (78,-19);
  \draw[->] (78,-28) -- node[msg] {3\ \texttt{ready}} (108,-28);
  \draw[->] (108,-37) -- node[msg, pos=0.6] {4\ \iflangja{データを読む}{read data}} (44,-37);
  \draw[->] (108,-46) -- node[msg] {5\ \texttt{msgsnd(ok)}} (78,-46);
  \draw[->] (78,-55) -- node[msg, pos=0.45] {6\ \texttt{ok}} (10,-55);
\end{tikzpicture}

  \caption{A message queue coordinates access to shared memory: P2 reads
    only after \enquote{ready}, and P1 continues only after \enquote{ok}.}
  \label{fig:l09-mq-sync}
\end{figure}

A semaphore (Lesson~10) is a synchronization construct of the operating
system that holds an integer value.  Used with the values 0 and~1 it acts
like a mutex: a process that finds the value~0 blocks; a process that finds
the value~1 decrements it to~0 and proceeds, and it increments the value
again when it leaves the critical section.  The System~V interfaces of the
three facilities follow one pattern (\cref{tab:l09-sysv-api}): a
\texttt{get} call creates or opens the object for a key, a \texttt{ctl} call
performs control operations including destruction with
\texttt{IPC\_RMID}, and the remaining calls do the actual work.

\begin{table}[htbp]
  \centering
  \caption{The System~V IPC interfaces.}
  \label{tab:l09-sysv-api}
  \small
  \begin{tabular}{@{}llll@{}}
    \toprule
    Facility & Create or open & Use & Control, destroy \\
    \midrule
    Shared memory & \texttt{shmget()} & \texttt{shmat()}, \texttt{shmdt()} & \texttt{shmctl()} \\
    Message queue & \texttt{msgget()} & \texttt{msgsnd()}, \texttt{msgrcv()} & \texttt{msgctl()} \\
    Semaphore set & \texttt{semget()} & \texttt{semop()} & \texttt{semctl()} \\
    \bottomrule
  \end{tabular}
\end{table}

\begin{keypoint}
  Shared memory needs the same synchronization as shared state in a
  multithreaded process.  Pthreads mutexes and condition variables work
  across processes when they are process-shared and reside in the shared
  region; message queues and semaphores provide the same coordination
  through the operating system.
\end{keypoint}

\section{IPC command-line tools}
\label{sec:l09-tools}

Because System~V IPC objects outlive the processes that created them,
Unix-like systems provide commands to inspect and remove them.\source{P3L3,
IPC command line tools; Kerrisk ch.~45.}  The command \texttt{ipcs} lists
all IPC facilities in the system---message queues, shared memory segments
and semaphore sets---with their keys, identifiers, owners and permissions;
\texttt{ipcs -m} restricts the list to shared memory.  The command
\texttt{ipcrm} removes a facility: \texttt{ipcrm -m} followed by a segment
identifier destroys that segment, which is the usual way to reclaim a
segment that a program failed to destroy.

\section{Designing communication over shared memory}
\label{sec:l09-design}

The operating system provides shared memory and then stays out of the
way.\source{P3L3, shared memory design considerations, how many segments,
design considerations.}  The choice of synchronization mechanism, the format
of the data and the protocol for passing them are left to the application.
This flexibility makes the programmer responsible for two design decisions
in particular.

\subsection{How many segments?}

Processes that communicate over shared memory can organize it in two ways
(\cref{fig:l09-segments}).
\begin{description}
  \item[One large segment] All communication uses a single segment.  The
    application then needs a memory manager that allocates portions of the
    segment to individual exchanges and frees them afterwards.
  \item[Many small segments] Each pair of communicating processes uses a
    segment of its own.  Because creating segments is costly and their
    number is limited, a pool of segments is created in advance, and the
    identifiers of the free ones are kept in a queue.  A process takes a
    segment from the pool and tells its peer the identifier through another
    mechanism, such as a message queue; after the exchange the segment
    returns to the pool.
\end{description}

\begin{figure}[htbp]
  \centering
  % Two ways of organizing shared memory among several processes.
% Usage: % Two ways of organizing shared memory among several processes.
% Usage: % Two ways of organizing shared memory among several processes.
% Usage: \input{figures/l09-segments}   (width ~118 mm)
\begin{tikzpicture}[x=1mm, y=1mm, font=\footnotesize\sffamily,
  blk/.style={draw, minimum height=8mm, inner sep=0pt, font=\scriptsize\sffamily},
  used/.style={blk, fill=ln-accent!20},
  freeb/.style={blk, fill=white},
  box/.style={draw, rounded corners=2pt, fill=ln-box, align=center,
              inner sep=1.5pt, font=\scriptsize\sffamily},
  cell/.style={draw, fill=white, minimum width=5mm, minimum height=5mm,
               inner sep=0pt, font=\scriptsize\sffamily},
  lab/.style={font=\scriptsize\sffamily, text=ln-muted, align=center},
  >={Stealth[length=1.6mm]}]
  % (a) one large segment managed by an allocator
  \node[used,  minimum width=12mm, anchor=west] (a1) at (0,0) {P1/P2};
  \node[freeb, minimum width=8mm,  anchor=west] (a2) at (a1.east) {};
  \node[used,  minimum width=12mm, anchor=west] (a3) at (a2.east) {P1/P3};
  \node[used,  minimum width=12mm, anchor=west] (a4) at (a3.east) {P2/P3};
  \node[freeb, minimum width=8mm,  anchor=west] (a5) at (a4.east) {};
  \node[box, text width=24mm] (alloc) at (26,18)
    {\iflangja{アロケータ\\（確保・解放）}{allocator\\(allocate / free)}};
  \foreach \t in {a1,a3,a4} \draw[->] (alloc.south) -- (\t.north);
  \node[lab] at (26,-9) {\iflangja{(a) 1つの大きなセグメント}{(a) one large segment}};
  % (b) pool of small segments
  \foreach \i/\s in {1/used,2/freeb,3/used,4/freeb} {
    \node[\s, minimum width=10mm] (s\i) at (58+13*\i,0) {seg \i};
  }
  \node[lab, anchor=east] at (86.5,18) {\iflangja{空き ID}{free ids}};
  \node[cell] (q2) at (90,18) {2};
  \node[cell] (q4) at (95,18) {4};
  \draw[->, dashed] (q2.south) -- (s2.north);
  \draw[->, dashed] (q4.south) -- (s4.north);
  \node[lab] at (90.5,-9) {\iflangja{(b) 小さなセグメントのプール}{(b) pool of small segments}};
\end{tikzpicture}
   (width ~118 mm)
\begin{tikzpicture}[x=1mm, y=1mm, font=\footnotesize\sffamily,
  blk/.style={draw, minimum height=8mm, inner sep=0pt, font=\scriptsize\sffamily},
  used/.style={blk, fill=ln-accent!20},
  freeb/.style={blk, fill=white},
  box/.style={draw, rounded corners=2pt, fill=ln-box, align=center,
              inner sep=1.5pt, font=\scriptsize\sffamily},
  cell/.style={draw, fill=white, minimum width=5mm, minimum height=5mm,
               inner sep=0pt, font=\scriptsize\sffamily},
  lab/.style={font=\scriptsize\sffamily, text=ln-muted, align=center},
  >={Stealth[length=1.6mm]}]
  % (a) one large segment managed by an allocator
  \node[used,  minimum width=12mm, anchor=west] (a1) at (0,0) {P1/P2};
  \node[freeb, minimum width=8mm,  anchor=west] (a2) at (a1.east) {};
  \node[used,  minimum width=12mm, anchor=west] (a3) at (a2.east) {P1/P3};
  \node[used,  minimum width=12mm, anchor=west] (a4) at (a3.east) {P2/P3};
  \node[freeb, minimum width=8mm,  anchor=west] (a5) at (a4.east) {};
  \node[box, text width=24mm] (alloc) at (26,18)
    {\iflangja{アロケータ\\（確保・解放）}{allocator\\(allocate / free)}};
  \foreach \t in {a1,a3,a4} \draw[->] (alloc.south) -- (\t.north);
  \node[lab] at (26,-9) {\iflangja{(a) 1つの大きなセグメント}{(a) one large segment}};
  % (b) pool of small segments
  \foreach \i/\s in {1/used,2/freeb,3/used,4/freeb} {
    \node[\s, minimum width=10mm] (s\i) at (58+13*\i,0) {seg \i};
  }
  \node[lab, anchor=east] at (86.5,18) {\iflangja{空き ID}{free ids}};
  \node[cell] (q2) at (90,18) {2};
  \node[cell] (q4) at (95,18) {4};
  \draw[->, dashed] (q2.south) -- (s2.north);
  \draw[->, dashed] (q4.south) -- (s4.north);
  \node[lab] at (90.5,-9) {\iflangja{(b) 小さなセグメントのプール}{(b) pool of small segments}};
\end{tikzpicture}
   (width ~118 mm)
\begin{tikzpicture}[x=1mm, y=1mm, font=\footnotesize\sffamily,
  blk/.style={draw, minimum height=8mm, inner sep=0pt, font=\scriptsize\sffamily},
  used/.style={blk, fill=ln-accent!20},
  freeb/.style={blk, fill=white},
  box/.style={draw, rounded corners=2pt, fill=ln-box, align=center,
              inner sep=1.5pt, font=\scriptsize\sffamily},
  cell/.style={draw, fill=white, minimum width=5mm, minimum height=5mm,
               inner sep=0pt, font=\scriptsize\sffamily},
  lab/.style={font=\scriptsize\sffamily, text=ln-muted, align=center},
  >={Stealth[length=1.6mm]}]
  % (a) one large segment managed by an allocator
  \node[used,  minimum width=12mm, anchor=west] (a1) at (0,0) {P1/P2};
  \node[freeb, minimum width=8mm,  anchor=west] (a2) at (a1.east) {};
  \node[used,  minimum width=12mm, anchor=west] (a3) at (a2.east) {P1/P3};
  \node[used,  minimum width=12mm, anchor=west] (a4) at (a3.east) {P2/P3};
  \node[freeb, minimum width=8mm,  anchor=west] (a5) at (a4.east) {};
  \node[box, text width=24mm] (alloc) at (26,18)
    {\iflangja{アロケータ\\（確保・解放）}{allocator\\(allocate / free)}};
  \foreach \t in {a1,a3,a4} \draw[->] (alloc.south) -- (\t.north);
  \node[lab] at (26,-9) {\iflangja{(a) 1つの大きなセグメント}{(a) one large segment}};
  % (b) pool of small segments
  \foreach \i/\s in {1/used,2/freeb,3/used,4/freeb} {
    \node[\s, minimum width=10mm] (s\i) at (58+13*\i,0) {seg \i};
  }
  \node[lab, anchor=east] at (86.5,18) {\iflangja{空き ID}{free ids}};
  \node[cell] (q2) at (90,18) {2};
  \node[cell] (q4) at (95,18) {4};
  \draw[->, dashed] (q2.south) -- (s2.north);
  \draw[->, dashed] (q4.south) -- (s4.north);
  \node[lab] at (90.5,-9) {\iflangja{(b) 小さなセグメントのプール}{(b) pool of small segments}};
\end{tikzpicture}

  \caption{Organizing shared memory for several pairs of processes:
    (a) one large segment whose portions an allocator hands out;
    (b) a pool of small segments whose free identifiers are kept in a queue.}
  \label{fig:l09-segments}
\end{figure}

\subsection{How large a segment?}

The second decision is the size of a segment, and what to do with data that
do not fit into one.
\begin{description}
  \item[Segment size equal to the data size] Every data item fits into one
    segment.  This works well when the data sizes are known and fixed, but
    it limits the size of the data to the maximum segment size.
  \item[Segment smaller than the data] Larger data are transferred in
    rounds, one portion at a time.  The processes then need a protocol that
    tracks the progress of the transfer, typically a data structure at the
    start of the segment that holds the synchronization constructs and
    records how much data the segment contains and whether the transfer is
    complete.
\end{description}

\begin{example}
  \label{ex:l09-rounds}
  A segment starts with the following header; the rest of the segment holds
  \SI{2}{\mega\byte} of payload.
  \begin{lstlisting}[language=C, numbers=none]
typedef struct {
  pthread_mutex_t mutex;  /* process-shared          */
  pthread_cond_t  cond;   /* process-shared          */
  size_t nbytes;          /* bytes in buf; 0 = empty */
  int    done;            /* set with last portion   */
  char   buf[];           /* payload                 */
} xfer_t;
  \end{lstlisting}
  Holding the mutex, the sender waits until \texttt{nbytes} is~0, copies
  the next portion into \texttt{buf}, sets \texttt{nbytes} (and
  \texttt{done} with the last portion), and signals.  The receiver waits
  until \texttt{nbytes} is non-zero, copies the portion out, resets
  \texttt{nbytes} to~0, and signals.  A transfer of \SI{5}{\mega\byte} takes
  $\lceil 5/2 \rceil = 3$ rounds:
  \begin{center}
    \small
    \begin{tabular}{@{}cllcc@{}}
      \toprule
      Step & Process & Action & \texttt{nbytes} (MB) & \texttt{done} \\
      \midrule
      1 & sender   & writes bytes 0--2\,MB & 2 & 0 \\
      2 & receiver & reads them            & 0 & 0 \\
      3 & sender   & writes bytes 2--4\,MB & 2 & 0 \\
      4 & receiver & reads them            & 0 & 0 \\
      5 & sender   & writes bytes 4--5\,MB & 1 & 1 \\
      6 & receiver & reads them, sees \texttt{done} & 0 & 1 \\
      \bottomrule
    \end{tabular}
  \end{center}
\end{example}

\begin{keypoint}[Lesson summary]
  IPC mechanisms move data across address spaces.  Message passing---pipes,
  message queues and sockets---sends every message through a kernel channel
  via ports, at the cost of system calls and copies for every message.
  Shared memory maps the same physical pages into several address spaces, so
  that after a one-time setup the processes communicate without the kernel;
  whether copying or mapping is cheaper depends on the data size and on how
  often the mapping is reused (\cref{eq:l09-copy-map}), a trade-off that
  Windows LPC exploits.  System~V segments (\texttt{shmget}, \texttt{shmat},
  \texttt{shmdt}, \texttt{shmctl}) and POSIX shared memory
  (\texttt{shm\_open}, \texttt{mmap}, \texttt{munmap}, \texttt{shm\_unlink})
  persist until explicitly destroyed; \texttt{ipcs} and \texttt{ipcrm}
  inspect and remove System~V objects.  Processes that share memory must
  synchronize, with process-shared Pthreads mutexes and condition variables
  placed in the region or with message queues and semaphores, and they must
  decide how many segments to use and how to transfer data larger than a
  segment.
\end{keypoint}

\end{document}
